% Generated draft: requires mathematical and source review.
\documentclass[11pt]{article}
\usepackage[margin=25mm]{geometry}
\usepackage{fontspec}
\setmainfont{Latin Modern Roman}
\usepackage{amsmath,amssymb,mathtools,unicode-math}
\setmathfont{Latin Modern Math}
\usepackage{enumitem,needspace,xurl,hyperref}
\setlength{\parindent}{0pt}
\setlength{\parskip}{5pt}
\setlength{\emergencystretch}{3em}
\newcommand{\N}{\mathbb{N}}
\newcommand{\Z}{\mathbb{Z}}
\newcommand{\Q}{\mathbb{Q}}
\newcommand{\R}{\mathbb{R}}
\newcommand{\C}{\mathbb{C}}
\newcommand{\F}{\mathbb{F}}
\newcommand{\A}{\mathbb{A}}
\newcommand{\PP}{\mathbb{P}}
\newcommand{\E}{\mathbb{E}}
\newcommand{\eps}{\varepsilon}
\newcommand{\dd}{\,\mathrm{d}}
\begin{document}
\section*{UnsolvedMath ID: 800725}
\textbf{Source problem code:} catalogue entry
\subsection{Definitions and mathematical statement}
Let \(F\) be a field, with no restriction on its characteristic. A central simple \(F\)-algebra is a finite-dimensional associative unital \(F\)-algebra whose center is \(F\) and whose only two-sided ideals are zero and the whole algebra. Two such algebras \(A\) and \(B\) are Brauer equivalent if \(M_r(A)\cong M_s(B)\) as \(F\)-algebras for some positive integers \(r,s\). The Brauer group \(\operatorname{Br}(F)\) consists of these equivalence classes, with multiplication induced by tensor product over \(F\) and identity represented by \(F\).

Determine the set
\[
\mathcal S=\bigl\{n\in\Z_{\geq 1}:\text{ there exists a field }F\text{ with }\lvert\operatorname{Br}(F)\rvert=n\bigr\}.
\]
Equivalently, for each positive integer \(n\), determine whether some field has exactly \(n\) Brauer classes. The source says ``natural numbers'' without specifying whether zero is included; the positive-integer formulation covers every possible finite group cardinality, since a group contains its identity.

The source motivates the question using the familiar examples \(\R\) and the iterated Laurent series fields \(\R((x_1))((x_2))\cdots((x_m))\), described there as examples with finite nonzero Brauer group. Here ``nonzero'' means nontrivial. This introductory qualification does not explicitly restrict the subsequent question to nontrivial Brauer groups, so \(n=1\) remains included. The reference to these familiar examples is motivational and supplies no restriction on the fields quantified over.
\subsection{Short English statement}
Which positive integers occur as the number of elements of the Brauer group of a field? The field may have any characteristic. The question includes the trivial Brauer group, although the source motivates it with finite nontrivial examples.
\subsection{Sources}
\begin{itemize}
\item Supplied local UnsolvedMath registry record: permanent ID \(800725\), exact source problem code \texttt{catalogue entry}; title: \emph{Open Problems on Central Simple Algebras --- Problem 40}. The record attributes the problem to Auel, Brussel, Garibaldi, and Vishne, with proposed year 2010.
\item Auel, Brussel, Garibaldi, and Vishne, \emph{Open Problems on Central Simple Algebras} (2010), Problem 40. Supplied source URL: \url{https://arxiv.org/abs/1006.3304}. The supplied provenance records access on 2026-10-16 and extraction by \texttt{arxiv\_latex}; these are supplied provenance claims, not independent source verification.
\item The local record labels the problem \emph{open} and records an update on 2026-10-16. This status is supplied curation, not an independently verified mathematical or literature assessment.
\end{itemize}
\end{document}
